\documentclass[11pt]{article}
\usepackage[margin=1.1in]{geometry}
\usepackage{amsmath, amssymb, amsthm}
\usepackage{booktabs}
\usepackage{microtype}
\usepackage[hidelinks]{hyperref}

\newtheorem{theorem}{Theorem}
\newtheorem{lemma}[theorem]{Lemma}
\newtheorem{proposition}[theorem]{Proposition}
\newtheorem{corollary}[theorem]{Corollary}
\theoremstyle{definition}
\newtheorem{question}[theorem]{Question}
\theoremstyle{remark}
\newtheorem{remark}[theorem]{Remark}
\newtheorem{example}[theorem]{Example}

\newcommand{\Q}{\mathbb{Q}}
\newcommand{\R}{\mathbb{R}}
\newcommand{\Gal}{\operatorname{Gal}}
\newcommand{\Kt}{\widetilde{K}}

\title{Minimizing the degree of a set of real algebraic numbers,\\
with an application: the snub cube has no constructible model}
\author{Mikhail Hogrefe\thanks{M.H.\ posed the problem (as Math StackExchange question
2953331), constructed its worked example --- backward, which is the easy direction only in the
sense that it took him and not us --- and asked the snub cube question that
Section~8 answers.}
\and Claude\thanks{Claude (Anthropic) worked out the analysis and wrote the note. Any errors
are the second author's.}}
\date{August 2026}

\begin{document}
\maketitle

\begin{abstract}
Given a finite set $S$ of real algebraic numbers, how small can the largest degree of an element
of $T(S)$ be made, over all affine maps $T(x) = ax + b$ with real algebraic coefficients and
$a \neq 0$? This question was asked on Math~StackExchange (question 2953331). We reduce it to a
problem about the \emph{invariant field} generated by the affine ratios of $S$, prove
group-theoretic lower bounds (the Galois group of the invariant field's closure must be a
section of a product of small symmetric groups), and give a linear-algebra upper-bound machine
that redoes the question's worked example in a four-line computation. For target degree $2$ the
two sides meet: we obtain a complete, finite decision procedure. As a further consequence of the
lower-bound method we prove the result that originally motivated the question: no set of points
similar to a snub cube has all coordinates expressible by quadratic irrationals --- indeed, none
has all coordinates in the field of constructible numbers, because the snub cube possesses a
similarity invariant of degree $3$, namely the square of the tribonacci constant.
\end{abstract}

\section{The problem}

Let $S = \{s_1, \dots, s_n\}$ be a finite set of real algebraic numbers. For an algebraic
number $w$, $\deg w = [\Q(w) : \Q]$, and $\deg S = \max_i \deg s_i$. We consider affine maps
$T(x) = ax + b$ with $a, b$ real algebraic and $a \neq 0$, and ask:
\[
d(S) \;=\; \min_T \, \deg T(S) \;=\; \min_{a, b} \, \max_i \, \deg(a s_i + b).
\]
Throughout, $(p(x))_k$ denotes the $k$-th real root of $p$ in increasing order, starting at
$k = 0$.

\begin{example}[the StackExchange example]\label{ex:mse}
Let
\[
S = \bigl\{ \alpha, \beta, \gamma \bigr\}
= \bigl\{ (8x^3+4x^2-10x-37)_0,\;
(x^6+4x^5-26x^4-90x^3+200x^2+100x-775)_1,\;
\]
\[
(x^6-2x^5-22x^4+22x^3+142x^2-272x-131)_1 \bigr\}
\approx \{1.74190,\; 4.40417,\; 4.88765\},
\]
so $\deg S = 6$. With $a = 1$ and $b = (x^3-2x^2+5)_0 \approx -1.24190$ one gets
$T(S) = \{1/2, \sqrt{10}, 1+\sqrt7\}$ and $\deg T(S) = 2$, which is optimal. The example was
constructed backward from its solution; the question asks for the forward direction.
\end{example}

An answer posted by user \emph{mathlove} exhibits, for this specific $S$, a four-parameter
family of optimal $(a,b)$, found by substituting the ansatz $a s_i + b = p_i + q_i \sqrt{r_i}$
into the three minimal polynomials and eliminating by hand. We will see below
(Section~\ref{sec:assessment}) that the family is precisely the orbit of the known solution
under rational affine reparametrization and conjugate selection, and that the computation,
while correct, is specific to the instance: the ansatz, the assumed degree of $b$, and the
branch choices in the elimination are all guided by the known answer. The purpose of this note
is to set up the general theory.

\section{The affine reduction and the invariant field}\label{sec:reduction}

Assume $n \ge 2$ and fix $s_1 \neq s_2$. An affine map is determined by the images
$w_1 = T(s_1)$, $w_2 = T(s_2)$; writing $u = w_1$ and $v = w_2 - w_1 \neq 0$, the remaining
images are
\begin{equation}\label{eq:affine}
w_i \;=\; u + v \lambda_i, \qquad
\lambda_i \;=\; \frac{s_i - s_1}{s_2 - s_1} \quad (i = 3, \dots, n),
\end{equation}
and conversely any real algebraic $(u, v)$ with $v \neq 0$ arises from exactly one $T$
(with $a = v/(s_2 - s_1)$, which is real and nonzero). We also set $\lambda_1 = 0$,
$\lambda_2 = 1$. The numbers $\lambda_3, \dots, \lambda_n$ are a complete system of invariants
of $S$ under affine maps. Define the \emph{invariant field}
\[
K \;=\; \Q(\lambda_3, \dots, \lambda_n),
\]
let $\Kt$ be its Galois closure, and $G = \Gal(\Kt/\Q)$.

\begin{proposition}\label{prop:bounds}
For every admissible $T$, the field $L_T = \Q(w_1, \dots, w_n)$ contains $K$. Consequently
$d(S) \ge 2$ unless all $\lambda_i$ are rational, in which case $d(S) = 1$. In the other
direction, $(u, v) = (0, 1)$ shows $d(S) \le \max_i \deg \lambda_i$.
\end{proposition}

\begin{proof}
$\lambda_i = (w_i - w_1)/(w_2 - w_1) \in L_T$. If all images were rational, all $\lambda_i$
would be rational; conversely if all $\lambda_i \in \Q$, the map with $(u,v) = (0,1)$ has
images $0, 1, \lambda_3, \dots, \lambda_n \in \Q$.
\end{proof}

$K$ is the field of moduli of the configuration; no affine change of coordinates can shrink
the joint field below it. The whole question is how the content of $K$ can be \emph{distributed}
across the individual images: in Example~\ref{ex:mse},
\begin{equation}\label{eq:lambda}
\lambda_3
= \frac{\gamma - \alpha}{\beta - \alpha}
= \frac{\tfrac12 + \sqrt7}{\sqrt{10} - \tfrac12}
= \frac{1 + 2\sqrt{10} + 2\sqrt7 + 4\sqrt{70}}{39}
\approx 1.18160,
\end{equation}
a primitive element of the biquadratic field $K = \Q(\sqrt7, \sqrt{10})$ of degree $4$; yet
the optimum spreads $K$ over three images of degree at most $2$.

\section{Galois-theoretic lower bounds}\label{sec:lower}

\begin{theorem}[section obstruction]\label{thm:section}
If some admissible $T$ achieves $\deg(a s_i + b) \le d_i$ for all $i$, then $G = \Gal(\Kt/\Q)$
is a \emph{section} (a quotient of a subgroup) of $S_{d_1} \times \cdots \times S_{d_n}$.
\end{theorem}

\begin{proof}
Let $N_i$ be the Galois closure of $\Q(w_i)$ and $\widetilde L$ the compositum of the $N_i$,
which is Galois over $\Q$. Restriction gives an injection
$\Gal(\widetilde L/\Q) \hookrightarrow \prod_i \Gal(N_i/\Q)$, since an automorphism trivial on
every $N_i$ is trivial on their compositum; and $\Gal(N_i/\Q)$ acts faithfully on the
$\deg w_i \le d_i$ roots of the minimal polynomial of $w_i$, so it embeds in $S_{d_i}$. By
Proposition~\ref{prop:bounds}, $K \subseteq L_T \subseteq \widetilde L$, hence
$\Kt \subseteq \widetilde L$ and $G$ is a quotient of $\Gal(\widetilde L/\Q)$.
\end{proof}

\begin{corollary}\label{cor:prime}
$d(S) \ge$ the largest prime divisor of $|G|$.
\end{corollary}

\begin{proof}
If $\max_i d_i = d$, then $|G|$ divides $|S_d|^n = (d!)^n$, so every prime divisor of $|G|$ is
at most $d$.
\end{proof}

Sharper obstructions come from the structure of sections of $\prod S_{d_i}$ with all
$d_i \le 3$: such a group $B \le \prod (S_3 \text{ or } S_2)$ has an abelian Sylow
$2$-subgroup of exponent $2$ (contained in a conjugate of $\prod C_2$), and a normal
$3$-subgroup $B \cap \prod C_3$ with elementary abelian $2$-quotient. Both properties pass to
quotients. This settles the smallest nontrivial case completely:

\begin{theorem}[quartic classification]\label{thm:quartic}
Let $n = 3$ and let $\lambda = \lambda_3$ have degree $4$.
\begin{itemize}
\item If $\Q(\lambda)$ is biquadratic (Galois group $V_4$), then $d(S) = 2$.
\item In all other cases ($C_4$, $D_4$, $A_4$, $S_4$), $d(S) = 4$.
\end{itemize}
In particular $d(S) = 3$ never occurs for a quartic invariant.
\end{theorem}

\begin{proof}
\emph{Lower bounds.} Suppose all three images have degree $\le 2$. Then
$L = \Q(w_1, w_2, w_3)$ is a compositum of at most three quadratic fields, hence a Galois
extension with group an elementary abelian $2$-group; $L \supseteq \Q(\lambda)$ forces
$L \supseteq \Kt$, so $G$ is a quotient of an exponent-$2$ abelian group. This eliminates
$C_4$ (exponent), and $D_4$, $A_4$, $S_4$ (nonabelian). Suppose instead all three images have
degree $\le 3$; by Theorem~\ref{thm:section}, $G$ is a section of a product of copies of
$S_3$ and $S_2$. For $G \in \{C_4, D_4, S_4\}$ the Sylow $2$-subgroup ($C_4$, $D_4$, $D_4$
respectively) is not a section of an elementary abelian $2$-group, either by exponent or by
nonabelianness. For $G = A_4$: the image of the normal $3$-subgroup of $B$ is a normal
$3$-subgroup of $A_4$, hence trivial ($A_4$'s proper normal subgroups are $1$ and the
$2$-group $V_4$), so $A_4$ would be a quotient of a subgroup of an elementary abelian
$2$-group --- absurd. Hence in the non-$V_4$ cases $d(S) \ge 4$, and $(u,v) = (0,1)$ attains
$4$.

\emph{Upper bound for $V_4$.} Let $F_1, F_2, F_3$ be the three quadratic subfields of the
biquadratic field $M = \Q(\lambda)$. The condition
\[
w_3 - w_1 = \lambda\,(w_2 - w_1), \qquad w_i \in F_i,
\]
is a system of $4$ $\Q$-linear equations (one per basis component of $M$) in
$2 + 2 + 2 = 6$ rational unknowns, so its solution space has dimension $\ge 2$. A solution
with $w_2 = w_1$ forces $w_3 = w_1$ and $w_1 = w_2 \in F_1 \cap F_2 = \Q$, so the degenerate
solutions form only the $1$-dimensional line $w_1 = w_2 = w_3 \in \Q$. Any nondegenerate
solution yields $T$ with all images of degree $\le 2$; since $\lambda \notin \Q$, at least
one image is irrational, and $d(S) = 2$ exactly.
\end{proof}

\begin{remark}\label{rem:a4}
The $A_4$ case shows that the section theorem is strictly stronger than naive degree
bookkeeping: an $A_4$-quartic field has no quadratic subfield, so \emph{every element} of
$\Q(\lambda)$ has degree $1$ or $4$ --- all powers of two --- yet degree $2$ (and even $3$) is
unreachable. Divisibility of degrees alone would not detect this.
\end{remark}

\section{The upper-bound machine, and Example~\ref{ex:mse} in four lines}\label{sec:upper}

Let $M$ be the Galois closure of $K$ and $G = \Gal(M/\Q)$ (assume for the moment that images
may be sought in subfields of $M$; Section~\ref{sec:descent} discusses when this loses
nothing). Subfields $F \subseteq M$ with $[F : \Q] \le d$ correspond to subgroups of $G$ of
index $\le d$. The machine, for a target $d$:

\begin{enumerate}
\item Enumerate assignments $(F_1, \dots, F_n)$ of such subfields, one per image.
\item Solve the $\Q$-linear system
$\{\, w_i - w_1 - \lambda_i (w_2 - w_1) = 0 \;(i = 3..n),\ w_i \in F_i \,\}$,
adding, if $M \not\subseteq \R$, the linear constraint that $w_1$ and $w_2 - w_1$ be fixed by
complex conjugation.
\item Any solution with $w_2 \neq w_1$ certifies $d(S) \le d$.
\end{enumerate}

For Example~\ref{ex:mse}: $M = K = \Q(\sqrt7, \sqrt{10})$, and we try
$F_1 = \Q$, $F_2 = \Q(\sqrt{10})$, $F_3 = \Q(\sqrt7)$, i.e.\
$w_1 = a$, $w_2 = b_0 + b_1\sqrt{10}$, $w_3 = c_0 + c_1\sqrt7$. With $\beta = b_0 - a$ and
$\lambda$ from~\eqref{eq:lambda}, expanding $w_3 - w_1 = \lambda (w_2 - w_1)$ over the basis
$\{1, \sqrt{10}, \sqrt7, \sqrt{70}\}$ gives
\[
\sqrt{10}:\; 2\beta + b_1 = 0, \qquad
\sqrt{70}:\; 4\beta + 2b_1 = 0, \qquad
1:\; c_0 - a = \tfrac{1}{39}(\beta + 20 b_1), \qquad
\sqrt{7}:\; c_1 = \tfrac{1}{39}(2\beta + 40 b_1).
\]
The $\sqrt{70}$ equation is a multiple of the $\sqrt{10}$ equation (a special feature of this
$\lambda$, which is what allows a \emph{rational} first image), leaving
$b_1 = -2\beta$, $c_0 = a - \beta$, $c_1 = -2\beta$ with $a, \beta$ free:
\[
(w_1, w_2, w_3) \;=\; \bigl(a,\; a + \beta - 2\beta\sqrt{10},\; a - \beta - 2\beta\sqrt7\bigr).
\]
At $(a, \beta) = (\tfrac12, -\tfrac12)$ this is $(\tfrac12, \sqrt{10}, 1 + \sqrt7)$ --- the
known optimum --- and the two free parameters are exactly composition with rational affine
maps. This replaces the sextic elimination of the posted answer by a $4 \times 6$ linear
system, and Theorem~\ref{thm:quartic} supplies the missing optimality proof.

\section{Degree 2 is completely decidable}\label{sec:d2}

For target degree $2$ the machine is \emph{complete}, with no assumption:

\begin{theorem}\label{thm:d2}
$d(S) \le 2$ if and only if all $\lambda_i$ lie in a single multiquadratic field (equivalently,
$K$ is multiquadratic; then $M = K$), and for some assignment of each $F_i \in \{\Q\} \cup
\{\text{quadratic subfields of } M\}$ the linear system of Section~\ref{sec:upper} has a
solution with $w_2 \neq w_1$. Both conditions are decidable by finite rational linear algebra.
\end{theorem}

\begin{proof}
($\Leftarrow$) is Section~\ref{sec:upper}; reality is automatic since $K \subseteq \R$ and a
real multiquadratic field is its own Galois closure. ($\Rightarrow$): suppose $(u, v)$,
$v \neq 0$, achieves images $w_i$ of degree $\le 2$. Then $L = \Q(w_1, \dots, w_n)$ is a
compositum of quadratic fields: a Galois extension $\Q(\sqrt{r_1}, \dots, \sqrt{r_k})$,
$r_j \in \Q$, with group $C_2^k$. All $\lambda_i \in L$, so $K$ and $M \subseteq L$ are
multiquadratic. It remains to move the images into subfields of $M$. Let
$H = \Gal(L/M) \le C_2^k$ and decompose $L = \bigoplus_\chi L_\chi$ into $H$-eigenspaces over
$M$; because $L$ is multiquadratic over $\Q$, each nonzero $L_\chi$ can be written
$M \cdot \sqrt{e_\chi}$ with a \emph{rational} representative $e_\chi$ (a product of the
$r_j$). Pick $\chi$ with the $\chi$-component of $v$ nonzero, and let $\psi : L \to M$ extract
the $\chi$-coordinate: $\psi(z) = z_\chi / \sqrt{e_\chi}$. Then $\psi$ is $M$-linear, so
applying it to~\eqref{eq:affine} gives a new admissible pair
$(\psi(u), \psi(v))$ with $\psi(v) \neq 0$ and images $\psi(w_i)$. Writing
$w_i = a_i + b_i \sqrt{d_i}$ with $a_i, b_i, d_i \in \Q$: if $\sqrt{d_i}$ has character
$\chi$, then $\psi(w_i) = b_i \sqrt{d_i}/\sqrt{e_\chi}$, an element of $M$ whose square is the
rational $b_i^2 d_i / e_\chi$, hence of degree $\le 2$; if not (and $\chi$ is nontrivial),
$\psi(w_i) = 0$; and for trivial $\chi$, $\psi(w_i) \in \Q(\sqrt{d_i}) \cap M$. In every case
the new images lie in quadratic-or-rational subfields of $M$.
\end{proof}

\section{The general algorithm and the descent question}\label{sec:descent}

For general $d$ the machine of Section~\ref{sec:upper} is a true algorithm --- finitely many
subgroup assignments, each a rational kernel computation --- and it always yields upper bounds,
while Theorem~\ref{thm:section} yields lower bounds. Its \emph{completeness} rests on:

\begin{question}[descent]\label{q:descent}
If some real algebraic $(u, v)$, $v \neq 0$, achieves $\deg(u + v\lambda_i) \le d_i$ for all
$i$, do there exist $(u', v') \in M^2$, $v' \neq 0$, achieving the same bounds?
\end{question}

We can prove this in two regimes beyond $d = 2$ (Theorem~\ref{thm:d2}):

\begin{lemma}[forced compositum]\label{lem:forced}
If $\prod_i d_i < 2\,[M : \Q]$, descent holds: $L = \Q(w_1, \dots, w_n)$ satisfies
$M \subseteq L$ and $[L : \Q] \le \prod_i d_i < 2\,[M:\Q]$, so $L = M$ already.
\end{lemma}

\begin{lemma}[trace descent]\label{lem:trace}
Let $E \supseteq M$ be Galois over $\Q$ containing $u, v$, let $H = \Gal(E/M)$, and
$S_i = \operatorname{Stab}(w_i)$ in $\Gal(E/\Q)$. If
$\operatorname{Tr}_{E/M}(c\,v) \neq 0$ for some $c$ in the fixed field of
$\langle S_1, \dots, S_n \rangle$, then descent holds: the pair
$\bigl(\operatorname{Tr}_{E/M}(c\,u), \operatorname{Tr}_{E/M}(c\,v)\bigr)$ lies in $M^2$ and,
because $H$ is normal and fixes every $\lambda_i$, its images
$\operatorname{Tr}_{E/M}(c\,w_i)$ are fixed by the image of $S_i$ in $\Gal(M/\Q)$, hence have
degree $\le [\,\Gal(E/\Q) : S_i\,] \le d_i$.
\end{lemma}

The gap is that the twisted trace can annihilate $v$; for junk extensions $L M / M$ of
exponent $2$ the eigenspace projection of Theorem~\ref{thm:d2} repairs this (quadratic
characters are rational), but for cubic junk only the plain trace is $\Q$-rationally
equivariant, and we do not know how to rule out its vanishing. In every worked example the
question is vacuous (Lemma~\ref{lem:forced} applies) or resolved, and we conjecture descent
holds in general; settling it would make the algorithm unconditionally complete. Independently
of descent, upper and lower bounds often meet, as in Theorem~\ref{thm:quartic}, and the
minimization of any monotone function of the degree profile (maximum, sum, lcm, \dots) uses the
same enumeration.

\section{Assessment of the posted answer}\label{sec:assessment}

In this framework, mathlove's computation is the $F$-assignment
$(\Q, \Q(\sqrt{10}), \Q(\sqrt7))$ instance of Section~\ref{sec:upper}, executed by polynomial
elimination instead of linear algebra. The ansatz (quadratic images, cubic $b$) and the branch
selections in the elimination are imported from the known solution rather than derived, other
branches of the case tree are not pursued, and optimality is not addressed. The four
parameters $(s, t, x, y)$ of the answer's family are the two rational-affine degrees of
freedom $(a, \beta)$ above together with two conjugate/root-index selections. So the answer is
a correct verification of the instance, not a general method --- the general method is the
subject of this note.

\section{Higher dimensions and the snub cube}\label{sec:snub}

For configurations $P = \{p_1, \dots, p_n\} \subset \R^m$ up to \emph{similarity} (the
higher-dimensional version of the original question), the role of the $\lambda_i$ is played by
the normalized Gram data
$\langle p_i - p_1, p_j - p_1\rangle / |p_2 - p_1|^2$, and the analogue of
Proposition~\ref{prop:bounds} holds with the same proof: every similarity invariant of this
kind lies in the field generated by the coordinates of \emph{any} similar copy. In particular:

\begin{proposition}\label{prop:ndlower}
If some ratio of squared distances of $P$ has degree that is not a power of $2$, then no
similar copy of $P$ has all coordinates in the field of constructible numbers. (Constructible
numbers lie in towers of quadratic extensions; finitely many of them generate a field of
$2$-power degree, and the invariant would be an element of it, so its degree would divide a
power of $2$.)
\end{proposition}

This applies to the snub cube. Let $t$ be the \emph{tribonacci constant}, the real root of
\[
t^3 = t^2 + t + 1, \qquad t \approx 1.83928675521416,
\]
irreducible over $\Q$ (no rational root), so $\deg t = 3$. A standard model of the snub cube
has as vertices the $12$ even permutations of $(1, 1/t, t)$ bearing an even number of minus
signs together with the $12$ odd permutations bearing an odd number (the mirror image swaps
the parities; since reflections are similarities, chirality plays no role below). Note
$1/t = t^2 - t - 1$, so all coordinates lie in the cubic field $\Q(t)$.

\begin{theorem}\label{thm:snub}
Let $P \subset \R^3$ be any set of points similar to the snub cube, and let $L$ be the field
generated by its coordinates. Then $3 \mid [L : \Q]$. Consequently the coordinates of a snub
cube can never all be constructible; a fortiori they cannot all be quadratic irrationals or
rationals, nor all lie in a multiquadratic field.
\end{theorem}

\begin{proof}
Work first in the standard model and put $s = 1/t$. Take the three vertices
\[
v = (1, s, t), \qquad w = (s, t, 1), \qquad x = (-s, -t, 1),
\]
(all lie in the even-permutation, even-sign class). Using $st = 1$, $s + t = t^2 - 1$, and
$s^2 + t^2 = (s+t)^2 - 2st = t^4 - 2t^2 - 1 = 2t$ (by $t^4 = 2t^2 + 2t + 1$):
\[
|v - w|^2 = (1-s)^2 + (s-t)^2 + (t-1)^2
= 2\bigl(1 + s^2 + t^2 - s - t - st\bigr)
= 2\,(1 + 2t - t^2),
\]
\[
|v - x|^2 = (1+s)^2 + (s+t)^2 + (t-1)^2
= 2\,(1 + st) + 2(s^2 + t^2) + 2s - 2t
= 4 + 2(s + t) = 2\,(t^2 + 1).
\]
Their ratio is
\[
\rho \;=\; \frac{|v-x|^2}{|v-w|^2} \;=\; \frac{t^2 + 1}{1 + 2t - t^2} \;=\; t^2,
\]
since $t^2 (1 + 2t - t^2) = t^2 + 2t^3 - t^4 = t^2 + 2(t^2+t+1) - (2t^2+2t+1) = t^2 + 1$.
(Numerically, $\rho = 8.76595\ldots/2.59120\ldots = 3.38298\ldots = t^2$.) Now $t^2 \notin \Q$
(else $\deg t \le 2$), and $\Q(t^2) \subseteq \Q(t)$ has degree dividing the prime $3$, so
$\deg \rho = 3$.

If $P$ is similar to the snub cube, the similarity carries $v, w, x$ to vertices
$v', w', x'$ of $P$ and preserves ratios of squared distances, so
$\rho = |v'-x'|^2 / |v'-w'|^2 \in L$. Hence $\Q(\rho) \subseteq L$ and
$3 = [\Q(\rho):\Q]$ divides $[L:\Q]$. A field generated by finitely many constructible
numbers has degree $2^k$, and $3 \nmid 2^k$.
\end{proof}

\begin{remark}
The same argument, with the same three-line computation, shows the snub cube cannot be
realized with coordinates in \emph{any} field of $2$-power degree. It is the polyhedral
analogue of the non-constructibility of the regular heptagon. By contrast, the regular
icosahedron and dodecahedron have models with every coordinate in $\Q(\sqrt5)$, and so are
constructible. The snub dodecahedron is amenable to the same method: its coordinate field
is a cubic extension of $\Q(\sqrt5)$, and the analogous distance-ratio computation produces an
invariant of degree divisible by $3$.
\end{remark}

\begin{remark}[limits of the method in higher dimensions]
The lower-bound side transfers to $\R^m$ wholesale, but the upper-bound side does not: rational
similarity invariants do \emph{not} guarantee rational coordinates. An equilateral triangle
has all invariants rational, admits rational coordinates in $\R^3$ (the unit basis vectors),
but none in $\R^2$ (a lattice triangle has rational area, while an equilateral triangle with
rational side lengths squared has area $\sqrt3 \cdot (\text{rational})$). Realizing a rational
Gram matrix by low-degree coordinates is a question about rational congruence of quadratic
forms (Hasse--Minkowski), and by M\"nev-type universality, realization spaces of polytopes can
force invariant fields of arbitrarily high degree, so hard instances are unavoidable. The
$1$-dimensional theory of Sections~\ref{sec:reduction}--\ref{sec:d2} is the part that admits a
clean algorithm.
\end{remark}

\section{Implementation notes}\label{sec:impl}

The question's $(p(x))_k$ notation is exactly the representation used by FLINT's
\texttt{qqbar} module (minimal polynomial plus isolated root enclosure), which Malachite plans
to port. The $1$-dimensional algorithm decomposes into:

\begin{itemize}
\item \textbf{Invariants:} $\lambda_i$ by \texttt{qqbar} arithmetic (resultant-based field
operations with enclosure refinement).
\item \textbf{Fields:} a primitive element for $K$; minimal polynomials require univariate
factorization over $\Q$ (van Hoeij), which is also the engine for the subfield algorithm of
van Hoeij--Kl\"uners--Novocin, which enumerates all subfields of a number field in polynomial
time --- precisely the $F_i$ enumeration for small $d$. Galois closures and $\Gal$ for the
lower bounds are available in Pari/GP for moderate degrees.
\item \textbf{Linear algebra:} the systems of Section~\ref{sec:upper} are exact rational
kernel computations; Malachite's \texttt{Rational} matrices suffice today.
\item \textbf{Degree 2:} Theorem~\ref{thm:d2} needs no Galois machinery at all --- only
square-class bookkeeping in a multiquadratic field --- and would make a good first target: it
already decides questions like ``does this configuration have a model over quadratic
irrationals?'', the question answered negatively for the snub cube above.
\end{itemize}

\medskip
\noindent\textbf{Source.} The original question is
\href{https://math.stackexchange.com/questions/2953331}{Math StackExchange 2953331,
``Minimizing the degree of a set of real algebraic numbers''} (2018); the posted answer
discussed in Section~\ref{sec:assessment} is by user mathlove.

\end{document}
